% Generated draft: requires mathematical and source review.
\documentclass[11pt]{article}
\usepackage[margin=25mm]{geometry}
\usepackage{fontspec}
\setmainfont{Latin Modern Roman}
\usepackage{amsmath,amssymb,mathtools,unicode-math}
\setmathfont{Latin Modern Math}
\usepackage{enumitem,needspace,xurl,hyperref}
\setlength{\parindent}{0pt}
\setlength{\parskip}{5pt}
\setlength{\emergencystretch}{3em}
\newcommand{\N}{\mathbb{N}}
\newcommand{\Z}{\mathbb{Z}}
\newcommand{\Q}{\mathbb{Q}}
\newcommand{\R}{\mathbb{R}}
\newcommand{\C}{\mathbb{C}}
\newcommand{\F}{\mathbb{F}}
\newcommand{\A}{\mathbb{A}}
\newcommand{\PP}{\mathbb{P}}
\newcommand{\E}{\mathbb{E}}
\newcommand{\eps}{\varepsilon}
\newcommand{\dd}{\,\mathrm{d}}
\begin{document}
\section*{UnsolvedMath ID: 2774}
\textbf{Source problem code:} catalogue entry
\subsection{Definitions and mathematical statement}
For genus $g\geq2$, let $\kappa=(k_1,\ldots,k_s)$ be a partition of $2g-2$ into positive integers. The stratum $\mathcal H_\kappa$ parametrizes pairs $(X,\omega)$ with $X$ a compact genus-$g$ Riemann surface and $\omega\ne0$ a holomorphic one-form whose zeros have exactly these multiplicities. One must specify whether zeros are labeled, whether the full nonzero differential or its projective class is retained, and whether the moduli orbifold, coarse space or a finite cover is used. These choices affect fundamental groups and monodromy, and the source intentionally asks about ``some version'' of the map.

Away from its zeros, integration of $\omega$ gives translation charts with transition maps $z\mapsto z+c$. To include zeros of order $k_i$, the translation-surface description must allow the corresponding singular points; a nonsingular translation atlas on the entire closed surface would not describe general strata. The source identifies the full space of differentials with the Hodge bundle, with fiber $H^{1,0}(X)$, and states complex dimension $2g+s-1$ for these positive-zero strata. It notes quasiprojectivity and finite CW homotopy type after a finite cover removing orbifold issues. The genus-one zero-free case needs a separate convention and is not silently folded into the positive-partition dimension formula.

The principal partition is $(1,\ldots,1)$ and the minimal one is $(2g-2)$. For each connected component $\mathcal H$ in the chosen moduli convention, the target program is to determine its homotopy type, decide whether it is a $K(\pi,1)$, compute its fundamental group, and decide injectivity of the appropriate monodromy
\[
\rho:\pi_1(\mathcal H)\longrightarrow\operatorname{Mod}(S_g).
\]
Here $\operatorname{Mod}(S_g)$ is a mapping class group of the underlying genus-$g$ surface with whatever markings the chosen version retains; in an orbifold model one uses the corresponding orbifold fundamental group. A $K(\pi,1)$ is a connected space with contractible universal cover, equivalently, in the usual CW setting, higher homotopy groups zero. The source asks for analogous topological comparisons as zeros merge from principal toward minimal strata, without defining a particular comparison map or a quantitative meaning of predictable change.

The background reports Kontsevich--Zorich [KZ03]'s classification of components, between one and three per stratum, and special hyperelliptic components for partitions $(2g-2)$ and $(g-1,g-1)$ related to configuration spaces in $\C$. It describes these as resolved cases and attributes the general asphericity/mapping-class-group proposal to [KZ97]. It reports Calderon--Salter [CS23]'s monodromy image for every nonhyperelliptic component in $g\geq5$ as a framed mapping class group, stabilizing the framing induced by translation charts on the punctured surface. This identifies an image, not automatically the full fundamental group: that identification additionally needs injectivity.

The notes report injectivity in the hyperelliptic setting but failure in genus three for $\mathcal H^{\mathrm{odd}}(4)$. They attribute to Looijenga--Mondello [LM14] a fundamental-group description $A(E_6)/Z(A(E_6))$, with Artin generators mapping to Dehn twists in an $E_6$ intersection pattern, and to Wajnryb [Waj99] a noncentral kernel element for the corresponding $A(E_6)\to\operatorname{Mod}(\Sigma_3)$ map. These descriptions are retained with the source's convention caveat; the central quotient in particular must not be transferred between full and projectivized or marked/unmarked strata without checking the cited definitions.

The literature triage, checked 2026-08-17, labels the general multipart program open and reports no complete classification. The primary registry status is ``open''. Detailed component/image results, selected aspherical cases and a specific noninjectivity example do not settle all homotopy, group and transition questions. No new theorem or convention-specific identification is claimed here.
\subsection{Short English statement}
Determine the topology, asphericity, fundamental groups and monodromy kernels of components of abelian-differential strata, and how these vary as zeros merge. The source records substantial component and monodromy information but no general solution; moduli, marking and scaling conventions are essential to precise group statements.
\subsection{Sources}
\begin{itemize}
\item UnsolvedMath local registry, permanent ID 2774, source code catalogue entry, ``Kirby Problem 2.26'', Kirby's Problems in Low-Dimensional Topology collection. Dataset attribution: Ulam AI and the UnsolvedMath Contributors; curation and metadata under CC BY 4.0.
\item \emph{K3: A New Problem List in Low-Dimensional Topology}, Problem 2.26 (2026), \url{https://aimath.org/pastworkshops/kirbylistrep.pdf}. The supplied literature triage records checks on 2026-08-17.
\item [KZ03] Maxim Kontsevich and Anton Zorich, \emph{Connected components of the moduli spaces of Abelian differentials with prescribed singularities}, Inventiones Mathematicae 153(3) (2003), 631--678, doi:10.1007/s00222-003-0303-x.
\item [KZ97] M. Kontsevich and A. Zorich, \emph{Lyapunov exponents and Hodge theory} (1997), arXiv:hep-th/9701164, as supplied without a direct URL.
\item [CS23] Aaron Calderon and Nick Salter, \emph{Framed mapping class groups and the monodromy of strata of abelian differentials}, Journal of the European Mathematical Society 25(12) (2023), 4719--4790, doi:10.4171/jems/1290.
\item [LM14] Eduard Looijenga and Gabriele Mondello, \emph{The fine structure of the moduli space of abelian differentials in genus 3}, Geometriae Dedicata 169 (2014), 109--128, doi:10.1007/s10711-013-9845-2.
\item [Waj99] Bronislaw Wajnryb, \emph{Artin groups and geometric monodromy}, Inventiones Mathematicae 138(3) (1999), 563--571, doi:10.1007/s002220050353.
\end{itemize}
\end{document}
